\section{Related Work}
The open intent detection methods can be categorized into three types: statistical methods, deep-learning methods and LLM-embedding methods.

\subsection{Statistical Methods for Intent Detection}
Statistical methods frame intent detection as closed-set classification, using hand-crafted features and classical decision rules.
\citet{wang2002combination} proposes a hybrid method that combines rule-based grammars with statistical classifiers, but the grammar must be hand-crafted for each domain.
To eliminate this dependence, \citet{haffner2003scaling} proposes a large-margin method using support vector machine (SVM) with discriminative training, showing that purely statistical classifiers can match hybrid systems.
However, this method still relies on manually designed features.
To automate feature extraction, \citet{tur2010what} proposes an n-gram method based on word co-occurrence statistics, though sparsity limits its effectiveness in open vocabularies.
To mitigate sparsity, \citet{schuurmans2020intent} compares Naive Bayes, SVM, and long short-term memory (LSTM) classifiers with sparse and dense features, finding that dense word embeddings replace discrete counts with continuous representations.
Despite these efforts, all statistical methods operate on surface-level features and cannot capture compositional semantics or long-range dependencies.

\subsection{Deep-Learning Methods for Intent Detection}
Deep-learning methods replace hand-crafted features with learned semantic representations, and introduce explicit mechanisms for rejecting unknown intents.
\citet{kim2014convolutional} proposes a convolutional neural network (CNN) method that extracts local semantic patterns from word embeddings; \citet{liu2016attention} proposes an attention-based recurrent method that captures sequential dependencies.
However, these closed-set classifiers produce overconfident predictions for OOS queries because softmax normalization forces all inputs into the known label space~\citep{hendrycks2017baseline}.
To address this, \citet{hendrycks2017baseline} proposes using maximum softmax probability (MSP) as a rejection score, but MSP does not reliably reflect epistemic uncertainty.
\citet{lee2018simple} proposes a distance-based method using Mahalanobis distance between test embeddings and class centroids under a per-class Gaussian assumption.
However, such parametric assumptions are often violated, and scoring-based methods do not explicitly model the decision boundary.
To construct the boundary directly, \citet{shu2017doc} proposes deep open classification (DOC), which replaces softmax with one-vs-rest sigmoid classifiers and uses Gaussian-fitted per-class thresholds.
However, DOC boundaries are coupled with the classification objective.
To decouple boundary learning, \citet{zhang2021adaptive} proposes adaptive decision boundary (ADB), learning spherical boundaries that balance empirical risk and open-space risk in a post-processing step.
However, ADB assumes one centroid per class and fails when intent semantics span multiple clusters.
\citet{li2025multi} proposes multi-granularity boundaries via granular-ball computing, partitioning each class into sub-regions with distinct centroids and radii.

\subsection{LLM-Embedding Methods for Intent Detection}
LLM-embedding methods leverage the world knowledge and reasoning abilities encoded in large pretrained models, enabling generalization to intents unseen during training.
\citet{arora2024intent} proposes a hybrid method that routes uncertain predictions from a fine-tuned sentence transformer to an LLM via in-context learning and chain-of-thought prompting.
However, this requires LLM inference for every uncertain query.
To reduce invocations, \citet{zaera2025efficient} proposes UDRIL (uncertainty-driven large language models triggering), which uses reconstruction-error scoring to selectively invoke a LoRA (low-rank adaptation) fine-tuned LLM only for ambiguous queries.
However, routing decisions can become unstable under distribution drift.
To address few-shot settings where labeled data is scarce, \citet{chen2025collaborating} proposes FCSLM, which uses the LLM for data augmentation during pretraining of a lightweight model, then employs multi-round inference with LLM refinement.
All LLM-embedding methods face a trade-off between reasoning quality and deployment cost.
